\documentclass[11pt]{article}
\usepackage[utf8]{inputenc}
\usepackage{amsmath,amssymb}
\usepackage{hyperref}
\title{Scalable Photocatalysis Water Splitting for Large-Scale Settings}
\author{Anonymous}
\date{}
\begin{document}
\maketitle

\begin{abstract}
This paper studies Photocatalysis Water Splitting. Grounded in a real, citable literature \cite{ref1,ref2,ref3,ref4,ref5,ref6,ref7,ref8},
we propose a focused method and report \textbf{illustrative} experiments.
All references below are real and verifiable (OpenAlex/arXiv); the reported
numbers are simulated to illustrate intended behaviour.
\end{abstract}

\section{Introduction}
Photocatalysis Water Splitting is an active research area. This work targets a concrete gap identified from
the recent literature and contributes a method together with an evaluation protocol.

\section{Related Work (real literature)}
The following works, retrieved from public bibliographic APIs, frame our study:
\begin{itemize}
\item Linic et al. (2011) — "Plasmonic-metal nanostructures for efficient conversion of solar to chemical energy", Nature Materials (cited 4818).
\item Ni et al. (2005) — "A review and recent developments in photocatalytic water-splitting using TiO2 for hydrogen production", Renewable and Sustainable Energy Reviews (cited 4133).
\item Cao et al. (2015) — "Polymeric Photocatalysts Based on Graphitic Carbon Nitride", Advanced Materials (cited 3673).
\item Wen et al. (2016) — "A review on g-C 3 N 4 -based photocatalysts", Applied Surface Science (cited 2973).
\item Yang et al. (2013) — "Roles of Cocatalysts in Photocatalysis and Photoelectrocatalysis", Accounts of Chemical Research (cited 2819).
\item Wang et al. (2019) — "Particulate Photocatalysts for Light-Driven Water Splitting: Mechanisms, Challenges, and Design Strategies", Chemical Reviews (cited 2754).
\end{itemize}

\section{Method}
We decompose the problem across shards with a communication-efficient aggregation rule that keeps overhead sublinear in scale. We instantiate this idea for Photocatalysis Water Splitting and analyse its cost/quality trade-off.

\section{Experiments (Illustrative)}
\begin{center}
\begin{tabular}{lcc}
System & Primary Metric & Efficiency \\ \hline
Strong baseline & 0.812 & 1.00$\times$ \\
Ablation & 0.828 & 0.92$\times$ \\
\textbf{Ours} & \textbf{0.851} & \textbf{0.71$\times$}
\end{tabular}
\end{center}
\emph{Metrics simulated to illustrate the intended effect; not measured.}

\section{Conclusion}
We connected a real literature to a concrete method for Photocatalysis Water Splitting. Future work will
validate the illustrative results with real experiments.

\begin{thebibliography}{9}
\bibitem{ref1} Suljo Linic, Phillip Christopher, David Ingram. Plasmonic-metal nanostructures for efficient conversion of solar to chemical energy. \emph{Nature Materials}, 2011. DOI:10.1038/nmat3151
\bibitem{ref2} Meng Ni, Michael K.H. Leung, Dennis Y.C. Leung et al.. A review and recent developments in photocatalytic water-splitting using TiO2 for hydrogen production. \emph{Renewable and Sustainable Energy Reviews}, 2005. DOI:10.1016/j.rser.2005.01.009
\bibitem{ref3} Shaowen Cao, Jingxiang Low, Jiaguo Yu et al.. Polymeric Photocatalysts Based on Graphitic Carbon Nitride. \emph{Advanced Materials}, 2015. DOI:10.1002/adma.201500033
\bibitem{ref4} Jiuqing Wen, Jun Xie, Xiaobo Chen et al.. A review on g-C 3 N 4 -based photocatalysts. \emph{Applied Surface Science}, 2016. DOI:10.1016/j.apsusc.2016.07.030
\bibitem{ref5} Jin‐Hui Yang, Donge Wang, Hongxian Han et al.. Roles of Cocatalysts in Photocatalysis and Photoelectrocatalysis. \emph{Accounts of Chemical Research}, 2013. DOI:10.1021/ar300227e
\bibitem{ref6} Qian Wang, Kazunari Domen. Particulate Photocatalysts for Light-Driven Water Splitting: Mechanisms, Challenges, and Design Strategies. \emph{Chemical Reviews}, 2019. DOI:10.1021/acs.chemrev.9b00201
\bibitem{ref7} Kazuhiko Maeda, Kazunari Domen. Photocatalytic Water Splitting: Recent Progress and Future Challenges. \emph{The Journal of Physical Chemistry Letters}, 2010. DOI:10.1021/jz1007966
\bibitem{ref8} Tiziano Montini, Michele Melchionna, Matteo Monai et al.. Fundamentals and Catalytic Applications of CeO<sub>2</sub>-Based Materials. \emph{Chemical Reviews}, 2016. DOI:10.1021/acs.chemrev.5b00603
\end{thebibliography}
\end{document}
